% Papiers de Sophie Germain — Français 9118, Gallica views 61–80 (batch 4).
%
% Pass: Opus 5 (claude-opus-5), 11 September 2026 — first pass, unchecked
% against the views by a human.
%
% WHAT THE TWENTY VIEWS HOLD. Letters received, and two mathematical notes,
% in eight pieces:
%   v. 61–62 (ff. 56–57)  Le Gendre to Germain, Paris 23 July 1821, on the
%                         printed « Recherches » she had just sent him; v. 63
%                         carries, on the facing leaf, the address panel of
%                         that same folded sheet.
%   v. 63–64 (ff. 58–59)  Le Gendre's mathematical note headed « Pag. 155 »,
%                         against Euler's solution sin ½ω = 0 for the
%                         vibrating lamina, and against Germain's following
%                         him into it; ends with the general resolution of
%                         Euler's equation and the case λ = ⅓.
%   v. 65    (f. 60)      A note in another hand, signed « G. Libri » and
%                         dated 1824, on decomposing a prime 4k+1 into two
%                         squares by a congruence on the central binomial.
%   v. 66    (f. 61)      Navier to Germain, Paris 2 August 1821.
%   v. 68–69 (ff. 63–64)  Chantel, chef de bataillon, to General Pernety,
%                         Brunswick 27 9bre 1806: his call on Gauss on
%                         Germain's behalf, with Gauss's address.
%   v. 70–71 (ff. 65–66)  Pernety to Germain, Cosel près Breslau 23 Xbre 1806,
%                         forwarding that letter, from the siege lines.
%   v. 72–74 (ff. 67–68)  Poisson to Germain, Paris 15 January 1816, on the
%                         Commission's reproach; f. 67 is the address panel.
%   v. 74–75 (ff. 69–70)  An undated invitation to dinner, « ce 17 pluviôse »,
%                         from Cloître Notre-Dame n° 8; signature a paraph.
%   v. 76, 77, 79 (ff. 71–73)  Three letters of d'Ansse de Villoison, 12 and
%                         14 July 1802 and one undated, on the Greek and Latin
%                         verses he had addressed to her and then burnt.
%
% SKIPPED. v. 67, 78, 80 — blank versos, carrying nothing but the offset of
% the facing letter. The gaps in the numbering are the record.
%
% NOTATION. Le Gendre draws a bar over e for a negative exponent, set here as
% \bar{e}^{x} for e^{-x}; a bar over an expression is his vinculum, i.e. a
% pair of brackets, set here as \overline{2-\lambda}. He writes « cot. » with
% the stop, and his beta is drawn like a figure 6 — read as beta from the
% sequence alpha, beta, gamma and set as \beta throughout. Libri's congruence
% sign is set as \equiv and his « (mod. p) » kept as he writes it. Underlining
% in the originals — Latin quotations, a journal title, « ma parole d'honneur »
% — is set as emphasis, so that the underline of \uncertain{} keeps its one
% meaning.
\documentclass[11pt,a4paper]{article}
% The preamble shared by every transcription of the mathematicians' manuscripts in Gallica.
%
% It rests on one idea: **uncertainty compounds, it does not resolve.** A
% transcription that smooths over a doubtful word has destroyed the only thing
% it was made for — a reader must be able to tell, at every point, what was on
% the page and what was supplied. The apparatus macros below are therefore the
% critical apparatus, not decoration.
%
% The same commands are recognised by `scripts/render.mjs`, which produces the
% site's reading view, and by `scripts/tei.mjs`. A macro added here must be
% added there too, or rendering fails loudly — which is the wanted behaviour.
%
% Comments here are English, like the rest of the repository. The documents
% this preamble sets are French, and that is deliberate: see the skills.

\ProvidesPackage{germain}[2026/09/15 Transcriptions of mathematicians' manuscripts in Gallica]

\usepackage{amsmath,amssymb,amsthm}
\usepackage{mathrsfs}
% Kept from the parent preamble so the renderer's diagram support stays
% available; these manuscripts draw no commutative diagrams, and nothing here uses it.
\usepackage{tikz-cd}
\usepackage[english,french]{babel}
\usepackage{fontspec}

% --- Greek ------------------------------------------------------------------
% The text face stays fontspec's default, Latin Modern, which has no Greek:
% XeTeX drops every glyph it lacks with a line in the log and nothing on the
% page, and the Greek words the manuscripts quote vanished from the PDFs. So
% the Greek and Greek Extended blocks get a XeTeX character class of their own,
% and entering or leaving a run of it switches the family to CMU Serif — the
% same Computer Modern design, with polytonic Greek — and back. Only the
% family changes: a Greek word in \textit or \textbf keeps its shape and series.
% The transitions to and from every other class carry the tokens that class
% has towards an ordinary letter, so French spacing before « ; » or inside
% « » still holds after a Greek word. No grouping, so no brace or \endgroup
% can come out unbalanced; maths is left alone, since XeTeX inserts nothing
% there. Kept identical in `exercices.sty`.
\makeatletter
\newfontfamily\germainGreekfont{cmun}[
  Extension=.otf, NFSSFamily=germaingreek,
  UprightFont=*rm, ItalicFont=*ti, SlantedFont=*sl,
  BoldFont=*bx, BoldItalicFont=*bi, BoldSlantedFont=*bl]
\newXeTeXintercharclass\germain@greek
\def\germain@greekrange#1#2{%
  \@tempcnta=#1\relax
  \loop
    \XeTeXcharclass\@tempcnta=\germain@greek
    \ifnum\@tempcnta<#2\advance\@tempcnta\@ne\repeat}
\germain@greekrange{"0370}{"03FF}
\germain@greekrange{"1F00}{"1FFF}
\def\germain@greekprev{\rmdefault}
\def\germain@greekin{%
  \edef\germain@greekprev{\f@family}\fontfamily{germaingreek}\selectfont}
\def\germain@greekout{\fontfamily{\germain@greekprev}\selectfont}
\def\germain@greekjoin{%
  \XeTeXinterchartokenstate=\@ne
  \@tempcnta=\z@
  \loop
    \ifnum\@tempcnta=\germain@greek\else
      \XeTeXinterchartoks\germain@greek\@tempcnta=\expandafter{%
        \expandafter\germain@greekout\the\XeTeXinterchartoks\z@\@tempcnta}%
      \XeTeXinterchartoks\@tempcnta\germain@greek=\expandafter{%
        \the\XeTeXinterchartoks\@tempcnta\z@\germain@greekin}%
    \fi
    \ifnum\@tempcnta<4095\advance\@tempcnta\@ne\repeat}
% babel-french sets its own classes, and saves and restores the interchar
% state, each time a language is selected: join again after it does.
\addto\extrasfrench{\germain@greekjoin}
\addto\extrasenglish{\germain@greekjoin}
\AtBeginDocument{\germain@greekjoin}
\makeatother

\usepackage{geometry}
\geometry{a4paper,margin=25mm,marginparwidth=28mm}
\usepackage{marginnote}
\usepackage{xcolor}
\usepackage[normalem]{ulem}
\setlength{\emergencystretch}{3em}
\usepackage{hyperref}
\hypersetup{colorlinks=true,linkcolor=black,urlcolor=black,pdfborder={0 0 0}}

% --- Batch metadata ---------------------------------------------------------
% Taken as they stand by the rendering script, which shows them at the head of
% the reading view. They fix what the file claims to cover. The macro names are
% the parent project's, so that the scripts stay shared; \folder here means the
% volume's slug — `fr-9115` — and \pages counts Gallica views.

\newcommand{\folder}[1]{\gdef\grothFolder{#1}}
\newcommand{\batch}[1]{\gdef\grothBatch{#1}}
\newcommand{\pages}[2]{\gdef\grothFirst{#1}\gdef\grothLast{#2}}
\newcommand{\foldertitle}[1]{\gdef\grothFoldertitle{#1}}

% The holder's own shelfmark, printed as they write it: « Français 9115 ».
\newcommand{\shelfmark}[1]{\gdef\grothShelfmark{#1}}

% The dating, copied verbatim from the holder's catalogue — never inferred
% here. « XIXe siècle » is the BnF's; a pass that dates a leaf from its content
% says so in a \note{}, as its own inference.
\newcommand{\dating}[1]{\gdef\grothDating{#1}}

% The status notice, printed in the title block and repeated as a diagonal
% watermark on every page of the PDF. The manuscripts are in the public
% domain; what this line declares is not a right but a quality — an
% unverified first pass by a machine. The skills write the line; a file
% without it gets no watermark, so leaving it out is visible in the output.
\newcommand{\watermark}[1]{\gdef\grothWatermark{#1}}

\gdef\grothFolder{?}\gdef\grothBatch{}\gdef\grothFirst{?}\gdef\grothLast{?}%
\gdef\grothFoldertitle{}\gdef\grothShelfmark{}\gdef\grothDating{}\gdef\grothWatermark{}

% --- The résumé -------------------------------------------------------------
\newenvironment{resume}
  {\small\setlength{\parskip}{0.45em}}
  {\par\medskip\hrule\medskip}

% --- Keywords ---------------------------------------------------------------
% In English, closing the modernised reading's résumé: the modern vocabulary
% under which to search for what the leaves construct. The manifest extracts
% them to tag the volume — this line is the single source of the tags.
\newcommand{\keywords}[1]{\par\smallskip\noindent
  {\footnotesize\textsc{Keywords} --- \textit{#1}}\par}

% --- The critical apparatus -------------------------------------------------

% The Gallica view number, in the margin. This is the anchor that lets a
% disagreement over a formula be settled by looking at *one* image rather than
% a volume of 750. The site uses it to turn the facsimile as the transcription
% is scrolled. A view is one scanned image: usually one side of a leaf.
\newcommand{\page}[1]{%
  \marginnote{\footnotesize\color{gray!70}\textsf{v.\,#1}}%
  \label{page:#1}\ignorespaces}

% The BnF's pencilled foliation, where the leaf carries it and a pass can read
% it: « 198r ». This is what the literature cites — « ff. 198r–208v » — and
% the one line that turns a citation into a view. Recorded, never inferred:
% a folio a pass cannot read is not written.
\newcommand{\folio}[1]{%
  \marginnote{\footnotesize\color{gray!50}\textsf{f.\,#1}}\ignorespaces}

% The modernised reading's coarser counterpart to \page{}: which views a
% section is drawn from.
\newcommand{\pagerange}[2]{%
  \marginnote{\footnotesize\color{gray!70}\textsf{v.\,#1--#2}}%
  \ignorespaces}

% Illegible: never guessed.
\newcommand{\ill}{\textcolor{red!60!black}{[\,\ldots\,]}}

% A doubtful reading: the word is given, and flagged as uncertain.
\newcommand{\uncertain}[1]{\uline{#1}}

% An editorial addition — a missing letter, an implied word.
\newcommand{\add}[1]{\textcolor{blue!50!black}{[#1]}}

% Struck out by the writer. What was crossed out is part of the document.
\newcommand{\struck}[1]{\sout{#1}}

% The transcriber's note, on the state of the page or the mathematics.
\newcommand{\note}[1]{\par\smallskip{\small\color{gray!60!black}\textit{[#1]}}\par\smallskip}

% A marginal note of hers, distinct from the body of the text.
\newcommand{\marginal}[1]{\par\smallskip\noindent\hspace{1em}%
  {\small\color{gray!60!black}#1}\par\smallskip}

% --- Title ------------------------------------------------------------------
% The title block is French, like the documents themselves.

\AddToHook{shipout/background}{%
  \ifx\grothWatermark\empty\else
    \begin{tikzpicture}[remember picture, overlay]
      \node[rotate=52, scale=3.4, align=center, text=gray!18,
            font=\sffamily\bfseries]
        at (current page.center) {\grothWatermark};
    \end{tikzpicture}%
  \fi}

\AtBeginDocument{%
  \begin{center}
    {\large\bfseries Manuscrits de mathématiciens (Gallica) ---
      \ifx\grothShelfmark\empty\grothFolder\else\grothShelfmark\fi}\\[2pt]
    {\itshape\grothFoldertitle}\\[4pt]
    {\small vues \grothFirst--\grothLast
      \ifx\grothBatch\empty\else\ (lot \grothBatch)\fi}\\[2pt]
    \ifx\grothDating\empty\else
      {\small\color{gray!60!black}Datation du catalogue : \grothDating}\\[2pt]
    \fi
    \ifx\grothWatermark\empty\else
      {\small\bfseries\color{red!45!gray}\grothWatermark}\\[2pt]
    \fi
    {\footnotesize\color{gray!60!black}Transcription établie d'après les images IIIF de Gallica.
     Source gallica.bnf.fr / Bibliothèque nationale de France.\\
     Les lectures incertaines sont soulignées, les illisibles notées [\,\ldots\,],
     les ajouts entre crochets ; \textsf{v.} numérote les vues de Gallica,
     \textsf{f.} la foliotation au crayon de la BnF.}
  \end{center}
  \vspace{1em}}

\endinput


\folder{fr-9118}
\shelfmark{Français 9118}
\batch{4}
\pages{61}{80}
\dating{1801-1900}
\watermark{Lecture automatique\\non vérifiée}
\foldertitle{Correspondance de Sophie GERMAIN avec les mathématiciens et savants Cauchy, Delambre, Fourier, Gauss, Le Gendre, d'Ansse de Villoison, etc.}

\begin{document}

\section{Lettre de Le Gendre — Paris, 23 juillet 1821}

\page{61}\folio{56}
\note{à gauche de la vue, le feuillet en regard porte la suscription d'un pli : « À Mademoiselle / Mademoiselle Sophie Germain » ; aucun chiffre au crayon ne s'y lit}

Mademoiselle,

J'ai reçu mardi dernier le mémoire que vous avez bien voulu m'envoyer, avec
beau papier, belle Couverture, et une petite lettre fort obligeante, mais trop
modeste. Je vous fais mon compliment bien sincère d'avoir enfin triomphé de
votre répugnance à rendre publiques des recherches qui vous ont coûté tant de
travaux. J'espère que vous n'aurez pas lieu de vous repentir de votre courage,
et que cette première émission qui étoit la plus difficile, sera bientôt suivie
de plusieurs autres qui obtiendront sans doute l'estime et le suffrage des
Connoisseurs.

Je n'ai pu encore que parcourir les premières pages de votre mémoire, et vous
pensez bien que je suis loin de pouvoir porter un jugement sur cet ouvrage qui
est du nombre de ceux qu'on ne peut apprécier que par

\page{62}\folio{57}
une étude longue et approfondie ; \struck{\ill{}} Sans doute vous repousseriez
vous même un jugement qui ne seroit fondé que sur un examen superficiel.

J'ai trouvé votre avertissement très bien rédigé, il présente fort nettement
l'état de la question ; Vous \struck{\ill{}} votre opinion de la manière la
plus modeste, et si l'on avoit quelque chose à Vous reprocher, ce seroit les
complimens dont en quelque sorte vous accablez le Géomètre dont vous combattez
l'opinion. Puisse-t-il répondre dignement à cet assaut de civilité ; c'est ce
que je désire plus que je n'espère.
\note{le verbe biffé, sous la rature, se lit peut-être « proposez » ; rien ne le remplace, en sorte que la phrase reste sans verbe}

J'ai été fâché de ne pas voir dans l'errata le mot \emph{Campanarum} à la place
de \emph{Campanorum} qui malheureusement est répété trois ou quatre fois.

Aussitôt que nous pourrons faire un petit séjour à Paris, nous nous
empresserons d'avoir l'honneur de vous voir. Ma femme se porte assez bien
maintenant. Elle vous fait mille tendres complimens, et Vous me permettrez,
Mademoiselle, de joindre mes hommages respectueux

Le Gendre

Paris le 23 juillet 1821.

\page{63}
\note{à gauche de la vue 63, la face extérieure du même pli, portant la suscription :}

\begin{quote}
à Mademoiselle / Mademoiselle Sophie Germain / Maison de M. Geoffroy, Médecin /
Rue de Braque au mara\add{is} / Paris
\end{quote}

\section{Note de Le Gendre sur la page 155}

\folio{58}
Pag. 155. (Legendre)

L'équation $\sin \frac{1}{2}\omega = 0$ n'est pas une conséquence nécessaire de
l'équation à résoudre ; elle vient d'un facteur qui a été introduit par la
Multiplication et qui est étranger à la solution du problème.

En effet la 1\textsuperscript{re} forme de l'équation générale (pag. 154) étant
\[
  0 = 2 - 2e^{2\omega}
      - \left(e^{\lambda\omega} - e^{\omega(2-\lambda)}\right)
        \left(e^{\lambda\omega} - \bar{e}^{\lambda\omega}\right)
        \left(\cot.\lambda\omega + \cot.(1-\lambda)\omega\right)
\]
Si on y fait $\lambda = \frac{1}{2}$, elle devient
\[
  0 = 2 - 2e^{2\omega}
      - \left(e^{\frac{1}{2}\omega} - e^{\frac{3}{2}\omega}\right)
        \left(e^{\frac{1}{2}\omega} - \bar{e}^{\frac{1}{2}\omega}\right)
        \left(2\cot.\tfrac{1}{2}\omega\right)
\]
or celle-ci n'est pas satisfaite par la supposition $\sin\frac{1}{2}\omega = 0$ ;
elle le seroit seulement par la supposition \struck{\ill{}}
$\frac{1}{2}\omega = 0$, ou $\frac{1}{2}\omega =$ à un infiniment petit, cas
dont on fait abstraction.

La solution $\sin\frac{1}{2}\omega = 0$, ou $\frac{1}{2}\omega = k\pi$ est
d'ailleurs inadmissible puisqu'elle donne des valeurs infinies pour les
coefficiens $\gamma$ et $\gamma'$, $\delta'$, pag. 153 (et toujours en faisant
$\lambda = \frac{1}{2}$). C'est ce qu'auroit dû remarquer Euler, lorsqu'il dit
pag. 156, \emph{Multiplicemus omnes Coefficientes par Sin
$\frac{1}{2}\omega$} ; on peut bien multiplier l'ordonnée d'une Courbe par une
Constante, afin de rendre cette Courbe sensible par une Construction
\emph{géométrique}, mais on ne peut pas multiplier par zéro.
\note{« par » est bien ce qui est écrit au milieu de la citation latine}

Il n'y a donc que la seconde solution qui soit légitime, et quant à celle-ci je
ne vois rien à lui objecter.

Lorsque Mlle Sophie a voulu considérer le cas général, elle est ce me semble
tombée dans la même erreur qu'Euler, en faisant $\sin\lambda\omega = 0$.

\page{64}\folio{59}
Cette solution est illusoire, elle résulte d'un facteur donné mal à propos à
l'équation, et elle auroit comme dans le cas de $\lambda = \frac{1}{2}$,
l'inconvénient de rendre infinis les Coefficiens, $\gamma$, $\gamma'$,
$\delta'$, \&c de la courbe.

Au reste excepté la 1\textsuperscript{re} solution qui demande quelques
tâtonnemens pour avoir la valeur précise de $\omega$, il est facile de
\struck{\ill{}} résoudre généralement l'équation d'Euler, pag. 154, savoir
\[
  0 = 2 - 2e^{2\omega}
      - \left(e^{\lambda\omega} - e^{\overline{2-\lambda}\,\omega}\right)
        \left(e^{\lambda\omega} - \bar{e}^{\lambda\omega}\right)
        \left(\cot.\lambda\omega + \cot.\overline{1-\lambda}\,\omega\right)
\]
En effet si on a bien saisi l'esprit de la résolution des six cas principaux,
on verra que passé la 1\textsuperscript{re} solution, et quelquefois même dans
cette 1\textsuperscript{re} solution, la quantité $e^{\omega}$ devient si
grande, qu'on peut négliger en toute sûreté $\bar{e}^{\omega}$ par rapport à
$e^{\omega}$, de même que $\bar{e}^{\lambda\omega}$ par rapport à
$e^{\lambda\omega}$. D'après ce principe l'équation précédente se réduit à
celle-ci
\[
  0 = -2e^{2\omega} + e^{2\omega}
      \left(\cot.\lambda\omega + \cot.\overline{1-\lambda}\,\omega\right)
\]
ou simplement
\[
  \cot.\lambda\omega + \cot.(1-\lambda)\omega = 2
\]
or ayant fait $\cot.\lambda\omega = x$, $\cot.(1-\lambda)\omega = y$, on
trouvera aisément, suivant les différentes valeurs de $\lambda$, une équation
algébrique entre $x$ et $y$, laquelle combinée avec l'équation $x + y = 2$,
donnera un nombre déterminé de solutions, par exemple
\[
  \lambda\omega = \alpha, \qquad \lambda\omega = \beta, \qquad
  \lambda\omega = \gamma, \qquad \&c
\]
De ces solutions on formera ensuite les solutions générales
\[
  \lambda\omega = \alpha + k\pi, \qquad \lambda\omega = \beta + k\pi, \qquad
  \lambda\omega = \gamma + k\pi, \qquad \&c
\]
$k$ étant un nombre à volonté.

Ainsi il y aura pour $\omega$ autant de suites de valeurs, que l'équation en
$x$ aura de racines.

Soit par exemple $\lambda = \frac{1}{3}$, il faudra satisfaire à l'équation
\[
  2 = \cot \tfrac{1}{3}\omega + \cot. \tfrac{2}{3}\omega
\]
or si l'on fait $\cot\frac{1}{3}\omega = x$, on aura
$\cot.\frac{2}{3}\omega = \frac{-1+xx}{2x}$, d'où
$x + \frac{x^2-1}{2x} = 2$, ou
\[
  3x^2 - 1 = 4x
\]
\note{une expression biffée, illisible sous la rature, précède cette dernière équation : c'est, autant qu'on en juge, la reprise de la précédente}
Et enfin $x = \frac{2 \pm \sqrt{7}}{3}$ : appellons $\alpha$ et $\beta$ les deux
angles compris entre 0 et 180°, qui donnent
$\cot\alpha = \frac{2+\sqrt{7}}{3}$, $\cot\beta = \frac{2-\sqrt{7}}{3}$, et
nous aurons généralement
\[
  \tfrac{1}{3}\omega = \alpha + k\pi
\]
\[
  \tfrac{1}{3}\omega = \beta + k\pi
\]
c'est-à-dire que les valeurs de $\omega$ formeront deux suites distinctes
\[
  3\alpha, \qquad 3\alpha + 3\pi, \qquad 3\alpha + 6\pi, \qquad \&c
\]
\[
  3\beta, \qquad 3\beta + 3\pi, \qquad 3\beta + 6\pi, \qquad \&c
\]
Chacune donnant lieu à une manière d'osciller de la lame.

Dans l'application il faudroit rechercher plus exactement les deux premiers
termes $3\alpha$, $3\beta$, mais les autres seront toujours suffisamment
approchés.

\section{Note signée G. Libri, 1824}

\page{65}\folio{60}
\note{d'une autre main que les pièces précédentes ; le feuillet en regard, à gauche de la vue, est blanc}

Si $p$ est un nombre premier de la forme $4k+1$, qui, à ce qu'on sait, peut se
mettre sous la forme $a^2 + b^2$, $a$ étant impair, $b$ pair : Je dis que l'on
aura
\[
  (k+1)(k+2)(k+3)\cdots 2k \equiv \pm\,2b, \quad (\mathrm{mod.}\; p)
\]
on aura aussi
\[
  \left[(k+1)(k+2)(k+3)\cdots 2k\right]^{2} \equiv \pm\,2b,
  \quad (\mathrm{mod.}\; p)
\]
et en conséquence
\[
  \frac{2}{1}\cdot\frac{6}{2}\cdot\frac{10}{3}\cdot\frac{14}{4}\cdot
  \frac{18}{5}\cdots\frac{4k-2}{k} \equiv \pm\,2a, \quad (\mathrm{mod.}\; p)
\]
il faut prendre le signe $+$ si $a$ est de la forme $4n+1$, et le signe $-$
lorsque $a$ est de la forme $4n-1$.
\note{les deux premières congruences portent le même second membre, bien que la seconde soit le carré de la première : la feuille est ainsi, et rien n'y est biffé}

On peut décomposer de cette manière directement le nombre $p$ en deux carrés.

G. Libri

1824

\section{Lettre de Navier — Paris, 2 août 1821}

\page{66}\folio{61}

Mademoiselle,

J'ai reçu avec reconnaissance l'ouvrage que vous avez bien voulu m'adresser. La
lecture que j'en ai faite m'a inspiré beaucoup d'intérêt, et j'apprécie autant
qu'il le mérite un écrit aussi remarquable, que bien peu d'hommes peuvent lire,
et qu'une seule femme pouvait faire.

J'ai l'honneur d'être avec respect, Mademoiselle, Votre très humble et très
obéissant serviteur

Navier

Paris, 2 août 1821.

\section{Lettre de Chantel au général Pernety — Brunswick, 27 9bre 1806}

\page{68}\folio{63}
Brunswick ce 27. 9\textsuperscript{bre} 1806.

\begin{quote}
A Monsieur Le Général pernetty chef de L'Etat major Général de L'Artillerie de
L'Armée
\end{quote}

Mon Général,

Appeine arrivée dans cette Ville, je me suis occupé de remplir votre
Commission, j'ai demandé a plusieurs personnes l'habitation de Monsieur Gauss
chez qui je fut pour prendre de ses nouvelles, de votre part, et de celles de
Mademlle sophie Germain, il me répondit ne pas avoir eu l'honneur de vous
Connoître ainsi que la Demoiselle, mais quil avait bien Connaissance de Madame
la lande a paris, \uncertain{aprec} avoir parlé des Différents articles contenu
dans votre instruction a moi remise, il me parut un peut Confus, et me chargea
de vous remercier infiniment des attentions que vous preniez a son Egard, je
lui demandai s'il voulait ecrire a Paris, de me remmettre la lettre que je vous
l'aurai faite tenir, donc vous chargiez de la faire rendre a sa destination, il
me repondit ni oui ni non sur cett'article, je sortis pour lors de chez lui en
le laissant avec madame son Epouse et son enfant, je fus chez Mr

\page{69}\folio{64}
\note{la suite occupe le verso du feuillet 63, à gauche de la vue ; le crayon « 64 » se lit sur le feuillet de droite, où la lettre s'achève}
Le Général de Div\textsuperscript{n} Buisson Gouverneur de cette ville pour le
reccommander, et surtout que j'avais eu l'honneur de connaitre Mr le Général
Buisson d'ancienne datte ; ce Général me répondit pour lors de faire tout pour
lui, en m'invitant a diner, avec Mr Gauss, Mr le Commandant de la place qui se
trouvait là dans ce moment me dit que cet homme lui avait déjà été
reccommandé par plusieurs personnes de mérite, je me suis licencié et je
retournai chez Mr Gauss pour le prier de vouloir bien venir diner avec moi chez
le Gouverneur, me l'aiant promis dans une heure d'ici je passerai le prendre et
nous irons ensemble — le fait est quil aura de Mr le Gouverneur et du
Commandant de la place toute l'estime et les douceurs qui seront a leur
pouvoir, chemin faisant je tacherai de lui parler affins quil vous écrive de la
manière donc je me suis acquitté de ma mission, et en même temps quil écrive a
paris s'il le juge, je lui laisse a cet effet votre adresse —

la sienne est a Monsieur le docteur Gauss logé chez Ritter steinweg N° 1917. a
Brunsvick, il jouis d'une bonne santé et me dit quil craignait un peut au
moment ou les troupes était ventrée, mais quil est resté a Brunsvick
tranquille, je l'ai rassuré, \uncertain{epprais} je ne doute pas que Messieurs
le Gouverneur, ainsij que le Command\textsuperscript{t} la place le rassureront
bien mieux sur cet article, j'ai couru la poste nuit et jour jusqu'à ce moment,
cette circonstance m'oblige a rester ici cet apprez midi, et demain matin de
bon heure je part pour me rendre a ma Destination —

Daignez agréer mon Général les sentiments du plus profond Respect avec le quel
j'ai l'honneur d'être.

Chantel chef de Bat\textsuperscript{n}

\section{Lettre du général Pernety — Cosel près Breslau, 23 Xbre 1806}

\page{70}\folio{65}
Cosel près Breslau 23 X\textsuperscript{bre} 1806

Mademoiselle,

Je ne puis mieux répondre aux demandes que votre amour pour les savants m'a
faites, qu'en vous envoyant la lettre de l'officier d'Artillerie que j'avois
chargé de savoir des nouvelles de M\textsuperscript{r} Gauss à Brunswick. Je
désire qu'elle satisfasse vos vœux pour cet émule d'Archimède, mieux traité que
lui comme vous le verrez. J'espère que cela me mettra dans le cas d'être chargé
quelquefois de vos intéressantes commissions. Je m'en acquitterai mieux
certainement que celles d'achats de chiffons des pays étrangers qu'à tort
parfois on me confie.

Me voici faisant un siège intendant et faisant gronder les, oui les, tonnerres,
brûlant des maisons des Églises, car les Clochers sont de bons points de mire
pour

\page{71}
\note{la fin de la lettre occupe le verso du feuillet 65 ; le crayon « 66 » se lit sur le feuillet blanc en regard}
les bombes, enfin faisant par Réflexion tout le mal que je peux à qui jamais ne
m'en fit aucun et que je ne connois pas, mais c'est le métier. On m'a accablé à
mon tour de Boulets, d'obus et de bombes et tout va le mieux du monde. Enfin
cet obstiné gouverneur de Breslau, prendra peut être un jour son parti, et il
fera bien pour la ville, et pour nous.

Je me flatte que votre santé est améliorée et que celles de vos Parents ainsi
que de Mademoiselle votre sœur se maintiennent en bon État ; tels sont du moins
les vœux les plus sincères de votre dévoué serviteur et admirateur

Pernety

\section{Lettre de Poisson — Paris, 15 janvier 1816}

\page{72}\folio{67}
\note{le feuillet ne porte que la suscription, écrite en travers ; le reste de la page n'est que le décalque de la lettre qui suit}

\begin{quote}
À Mademoiselle / Mademoiselle Germain / à Paris
\end{quote}

\page{73}\folio{68}
Paris, ce 15 janvier, 1816.

Mademoiselle,

M. Hallé vient de me remettre une lettre que vous me faites l'honneur de
m'adresser, et qui contient plusieurs questions relatives à votre mémoire. Le
reproche que la Commission lui a fait porte moins sur l'hypothèse dont vous
êtes partie que sur la manière dont vous avez appliqué le calcul à cette
hypothèse. Le résultat auquel ce calcul vous a conduit ne s'accorde avec
\uncertain{le mien}, que dans le seul cas où la surface s'écarte infiniment peu
d'un plan, soit
\note{lecture douteuse : le mot pourrait aussi bien se lire « l'ancien »}

\page{74}
dans l'état d'équilibre, soit dans l'état de mouvement. On imprime maintenant
mon mémoire, et je me propose de vous en offrir un exemplaire, aussitôt que
l'impression sera achevée. Permettez donc, Mademoiselle, que nous ajournions
ces discussions à l'époque où vous aurez pu comparer mes résultats aux vôtres.

Agréez l'hommage de mon respect et de ma haute considération.

Poisson

\section{Billet — Cloître Notre-Dame, ce 17 pluviôse}

\folio{69}
\note{sur un feuillet plus petit, monté à droite de la vue 74 ; d'une main que la pièce ne nomme pas}

Mademoiselle,

Duodi, c'est-à-dire Dimanche prochain il y aura chez moi un dîner, presque de
tous hommes. La majeure partie des convives ne vous est point étrangère. Vous
leur ferez grand plaisir et vous comblerez de bonheur le maître de la maison,
si vous vouliez bien être de la partie. Point de \uncertain{d. d.} puisque vous
ne vous êtes pas encor raccommodée avec lui. Vous trouverez dans ma Chartreuse
du beurre frais, des pommes de terre, des betteraves, des maches et quelque
autre aliment d'anachorète, qui n'irritera pas votre palais. Je compte sur une
de mes parentes, excellente femme, pour laquelle je vous demanderai indulgence
plénière,

\page{75}\folio{70}
car elle ne connoît de géométrie que les figures les plus naturelles. Vous en
jugerez par l'échantillon de son travail, que j'aurai l'honneur de vous
présenter. Comme j'ai bien à cœur de ne me point brouiller avec monsieur votre
père, vous pourriez lui promettre d'avance que vous serez reconduite chez lui
en sûreté, à l'heure qu'il désireroit. Quand vous verrez madame Lherbette, je
vous prie de l'assurer de mon respect, en lui disant que je la trouve une bien
bonne Commissionaire. Je désirerois réussir aussi bien qu'elle dans
\struck{celle} la commission, dont je m'acquitte aujourd'hui pour mon propre
compte.

Je suis avec respect, mademoiselle, Votre très humble serviteur \ill{}

Cloître Notre Dame, n° 8, ce 17 pluviôse
\note{la signature est un paraphe dont aucune lettre ne se laisse lire. Le millésime manque ; la lecture qui suit est celle de la présente transcription et n'est pas sur la feuille : « duodi » est le second jour de la décade, donc ici le 22 pluviôse, et le 22 pluviôse ne tombe un dimanche, dans ces années-là, que l'an VII — soit le 10 février 1799, le billet étant alors du 5 février}

\section{Trois lettres de d'Ansse de Villoison}

\page{76}\folio{71}
\note{à gauche de la vue, la face extérieure du pli : « À la Citoyenne Germain (Sophie), rue S\textsuperscript{te} Croix la Bretonnerie, près la rue de Moussy » — le double s y est écrit avec l's long}

Madame,

Je trouve en rentrant la lettre dont vous m'honorez et m'empresse de vous
donner sur le champ ma parole d'honneur que vos ordres et ceux de Mademoiselle
votre fille sont déjà ponctuellement exécutés, que j'ai brûlé ma pièce de vers
grecs, et voudrois pouvoir anéantir de même les Latins ; que je me contenterai
d'admirer désormais Mademoiselle votre fille dans le plus respectueux silence,
et de vous regarder comme la plus heureuse des mères et la plus digne d'envie.

J'oserai prendre la liberté de prier Mademoiselle Sophie d'agréer un Exemplaire
de la seconde Édition de Paris qui va paroitre sous très peu de jours, avec
l'addition que j'ai eu l'honneur de lui communiquer et qu'on m'annonce Madame
devoir être incessamment suivie d'une traduction en vers françois dont s'occupe
maintenant une Dame que je n'ai pas l'avantage de connoître. Je me reprocherai
toute ma vie d'avoir composé cette pièce, qui a pu blesser l'excessive modestie
de Mademoiselle votre fille. Je la supplie, ainsi que Mademoiselle sa sœur, de
vouloir bien recevoir mes excuses, et les assurances du vif et éternel regret
et du respect avec lequel je suis, Madame,

Votre très-humble et très-obéissant serviteur d'Ansse de Villoison

Ce lundi soir à minuit, 12 juillet 1802

\page{77}
\note{aucun chiffre au crayon ne se lit sur cette vue. Au verso du feuillet 71, à gauche, près du pli, une colonne de chiffres au crayon, très effacée et coupée par la pliure : \ill{}}

Mademoiselle,

J'ose prendre la liberté de vous offrir cy-joint un Exemplaire de la nouvelle
Édition de ma malheureuse pièce, avec les corrections et additions que je vous
avois annoncées. M\textsuperscript{r} Pougens, Mademoiselle, l'avoit insérée
dans le troisième numéro de la troisième année de sa \emph{Bibliothèque
Françoise}, avant que je pusse soupçonner que l'hommage de la vérité
choqueroit votre modestie aussi rare que vos talens. Je vous réitère avec mes
excuses, et \struck{de} l'expression de mes vifs et éternels regrets,
\emph{ma parole d'honneur} que jamais je ne me permettrai de parler de vous,
Mademoiselle, dans aucun écrit, et que mon admiration sera toujours muette, et
enchaînée par le désir d'obtenir mon pardon d'une erreur, ou d'une faute
involontaire, et par le profond respect que j'ai voué à Madame votre mère et à
Mademoiselle votre sœur, et avec lequel j'ai l'honneur d'être Mademoiselle

Votre très-humble et très-obéissant serviteur d'Ansse de Villoison

rue de Bièvre n° 22

Ce 14 juillet 1802.

P.S. Vous m'avouerez, Mademoiselle, que si vous êtes la seule demoiselle qui
possède si supérieurement les mathématiques, vous êtes aussi la seule qui ait
couru et redouté le danger d'un poème grec. En conscience, j'en appelle à
Mademoiselle votre sœur, qui est si bonne, ne pourriez-vous pas m'accorder ma
grâce, ne fut-ce que pour la singularité du fait ?

\page{79}\folio{73}

Mademoiselle,

Malgré la proscription fatale dont vous avez frappé un célèbre Astronome et ses
amis, je n'ai pas pu me dispenser de rendre hommage à la vérité, et je me suis
empressé de vous offrir les prémices d'une pièce de vers Latins de ma
composition, qui va paroître ces jours-cy dans le \emph{Magasin
Encyclopédique}, vous y verrez p. 239 une foible partie de la justice que je
vous rends, Mademoiselle, et qui vous est due à tant de titres.

Je serois trop heureux si vous vouliez présenter à Mademoiselle votre sœur, et
agréer l'hommage de l'admiration, et du respect avec lequel je suis
Mademoiselle

Votre très-humble et très-obéissant serviteur d'Ansse de Villoison

rue de Bièvre, n° 22.

\end{document}
